% TOP_PROBLEM: {"id": 3245, "title": "List chromatic number and maximum degree of bipartite graphs", "category_id": 3, "category": {"id": 3, "name": "graph_theory", "display_name": "Graph Theory", "description": "Problems involving graphs, networks, and their properties.", "slug": "graph-theory", "order_index": 3, "created_at": "2026-07-31T15:26:25.671Z"}, "status": "open", "problem_number": "OPG-47358", "set_id": 12, "statusQualification": "The logarithmic statement is interpreted for maximum degree at least two; an equivalent all-degree formulation is supplied."}
% FIRST_POSTED: 2026-10-02
% ATTEMPT_STATUS: unresolved
% UnsolvedMath ID 3245; source catalogue code OPG-47358
% !TeX program = lualatex
% Research attempt by GPT-6 Astra Ultra; no claim of a new complete solution.
\documentclass[11pt,a4paper]{article}
\usepackage[margin=25mm]{geometry}
\usepackage{fontspec}
\setmainfont{lmroman10-regular.otf}[BoldFont=lmroman10-bold.otf,
 ItalicFont=lmroman10-italic.otf,BoldItalicFont=lmroman10-bolditalic.otf]
\usepackage{amsmath,amssymb,mathtools}
\usepackage{unicode-math}
\setmathfont{latinmodern-math.otf}
\AtBeginDocument{\let\setminus\smallsetminus}
\usepackage{xurl,hyperref}
\hypersetup{colorlinks=true,urlcolor=blue}
\setcounter{secnumdepth}{0}
\setlength{\parindent}{0pt}
\setlength{\parskip}{5pt}
\setlength{\emergencystretch}{3em}
\begin{document}
\section*{List chromatic number and maximum degree of bipartite graphs}\label{3245}
{\small UnsolvedMath ID 3245 · OPG-47358 · Graph Theory · OpenGarden}
\subsection{Definitions and mathematical statement}
Let $G=(V,E)$ be a finite simple bipartite graph: there is a partition
$V=A\sqcup B$ such that every edge has one end in each part. Let
$\Delta(G)$ be the largest vertex degree. A graph is $k$-choosable if, for
every assignment of a $k$-element set $L(v)$ of colours to each vertex,
there is a proper colouring $c$ with $c(v)\in L(v)$. The least such $k$
is the list chromatic number $\chi_\ell(G)$; for a nonempty edgeless graph
it equals one.

The logarithmic-bound conjecture asks whether there exists an absolute
constant $C>0$ such that, for every finite bipartite graph with
$\Delta(G)\ge2$,
\[
                      \chi_\ell(G)\le C\log\Delta(G).
\]
Here $\log$ denotes the natural logarithm. Changing its base merely
changes $C$. The restriction $\Delta\ge2$ makes the formula meaningful:
$\log1=0$ cannot bound the choice number of an edge. An equivalent
all-degree formulation is that some constant $C'>0$ satisfies
$\chi_\ell(G)\le C'\log(\Delta(G)+2)$ for every nonempty bipartite graph.

The same constant must work independently of the order, the relative
sizes of the two parts, and all overlaps among lists. Bipartite graphs
have ordinary chromatic number at most two, but that does not give a
two-colour list assignment compatible with arbitrary vertex lists.
The target concerns this additional list-colouring difficulty rather
than ordinary two-colourability.
\subsection{Short English statement}
Is the list chromatic number of every bipartite graph bounded by a constant times the logarithm of its maximum degree?
\subsection{Research attempt}
\paragraph{Scope and process.}
This is a GPT-6 Astra Ultra research attempt dated 2 October 2026, using
primary-source browsing and elementary mathematical checks. The canonical
definition above is retained. Results below are restricted results or
reductions, not a claim of a new solution to the entire catalogue question.
No formal proof assistant or external mathematical referee checked this text.


\paragraph{Origin, literature and plan.}
This is the Alon--Krivelevich logarithmic maximum-degree conjecture:
the original paper [R1] states it as Conjecture 5.1. It is distinct
from the older minimum-order problem for nonchoosable bipartite graphs.
Bradshaw, Mohar and Stacho [R2] give a general upper bound on the
scale $\Delta/\log\Delta$, which does not supply the logarithmic
scale requested here. We prove the target for components of order
polynomial in $\Delta$, give explicit logarithmic lower-bound examples,
and identify the extra dependence that blocks the simple probability
argument in unrestricted graphs.

\paragraph{A palette-splitting lemma with its quantifiers checked.}
Let $G$ be bipartite with classes $X,Y$ and $N$ vertices, and assign
a list $L(v)$ of $q$ colours to each vertex. Independently send every
colour name appearing in any list to one of two palettes, each with
probability $1/2$. A vertex of $X$ is bad if its list has no colour
in the first palette; a vertex of $Y$ is bad if its list has none
in the second. Each individual bad event has probability $2^{-q}$,
regardless of list overlaps. A union bound gives probability at most
$N2^{-q}$ that any vertex is bad.

Whenever $N2^{-q}<1$, there exists a split for which every vertex
has an allowed colour in its designated palette. Choose any such
colour at each vertex. Edges cross the bipartition and the palettes
are disjoint, so the colouring is proper. This proves
\[
 \chi_\ell(G)\le\lfloor\log_2N\rfloor+1.
\]
Strict inequality in the probability bound matters when $N$ is a
power of two; the displayed extra one handles that case.

Suppose every connected component has at most $\Delta^a$ vertices
for some fixed $a>0$ and $\Delta\ge2$. Apply the lemma separately
to components. Even if colour names occur in different components,
the palette choices may differ because there are no cross-component
edges. We obtain
\[
 \chi_\ell(G)\le a\log_2\Delta+1
 \le\frac{a+1}{\log2}\log\Delta.
\]
This verifies the conjectured form on that family, with a constant
depending only on the fixed exponent $a$. In particular it applies
to balanced complete bipartite graphs, whose order is twice their
maximum degree, by the original $N$-bound.

\paragraph{Explicit examples requiring logarithmically many list colours.}
Fix $q\ge2$, take a colour universe of size $2q-1$, and put
$n=\binom{2q-1}{q}$. On each side of $K_{n,n}$ assign every
$q$-element subset of that universe once as a list. In any colouring,
the colours used on one side must meet every $q$-subset. Such a
hitting set has at least $q$ elements, since a set of at most $q-1$
elements has a complement containing a $q$-subset. Both sides thus
use at least $q$ colours. The total universe has only $2q-1$, so
the two used-colour sets intersect. Every cross-pair is adjacent,
contradicting properness. Therefore $\chi_\ell(K_{n,n})>q$.

Since $n\le2^{2q-1}<4^q$, this family gives a lower bound of
logarithmic order in $\Delta=n$. It does not contradict the target;
it proves that a bounded constant independent of degree could not
replace its logarithm. For $q=2$ the construction is $K_{3,3}$
with lists $\{1,2\},\{1,3\},\{2,3\}$ on each side.
The companion script [R3] checks noncolourability for $q=2,3$ by
enumerating all partitions of the finite colour universe between the
sides, rather than assuming independent vertex choices.

\paragraph{Why the order dependence is a real gap in this proof.}
A connected bipartite graph can have fixed maximum degree and
arbitrarily many vertices: paths already do. Thus $\log N$ cannot
be substituted by $O(\log\Delta)$ in the unrestricted palette lemma.
Paths themselves are two-choosable by colouring from an endpoint, so
they are not counterexamples; they expose a weakness of the method.
Likewise a colour may appear in lists at arbitrarily many far-apart
vertices. The dependency graph of palette-failure events is therefore
not bounded by $\Delta$ merely because the graph has bounded degree.
Invoking a local lemma without an additional dependency analysis would
not remove the unwanted order factor.

\paragraph{Final status.}
\textbf{UNRESOLVED.} The order-controlled upper bound and explicit
lower-bound family are proved. The missing step is a colouring argument
whose cost depends on graph degree while tolerating arbitrarily large
components and unrestricted global overlaps among lists.

\subsection{Sources}
\begin{itemize}
\item[S1] ULAM.AI and the UnsolvedMath Contributors, supplied problems.json, record 3245 (OPG-47358), OpenGarden collection. Catalogue identity and source transcription; CC BY 4.0.
\url{https://huggingface.co/datasets/ulamai/UnsolvedMath}.
\item[S2] Open Problem Garden, List chromatic number and maximum degree of bipartite graphs. Original problem and discussion; the mathematical formulation below is expanded from the supplied source record.
\url{http://www.openproblemgarden.org/op/list_chromatic_number_and_maximum_degree_of_bipartite_graphs}.

\item[R1] Noga Alon and Michael Krivelevich, \emph{The choice number
of random bipartite graphs}, Annals of Combinatorics 2 (1998), 291--297.
Original Conjecture 5.1 checked 2 October 2026.
\url{https://web.math.princeton.edu/~nalon/PDFS/choiceb4.pdf}.
\item[R2] Peter Bradshaw, Bojan Mohar and Ladislav Stacho,
\emph{Bipartite graphs are $(\frac45-\varepsilon)\Delta/\log\Delta$-choosable},
2024 preprint, primary abstract checked 2 October 2026.
\url{https://arxiv.org/abs/2409.01513}.
\item[R3] GPT-6 Astra Ultra, exact batch certificates, 2 October 2026.
Repository script \texttt{scripts/check\_attempt\_batch03\_colouring\_2026\_10\_02.py}.

\end{itemize}
\end{document}
